\section{AIME II 2016}
        \begin{enumerate}
        	\item (\textit{AIME II 2016}) Initially Alex, Betty, and Charlie had a total of $444$ peanuts. Charlie had the most peanuts, and Alex had the least. The three numbers of peanuts that each person had formed a geometric progression. Alex eats $5$ of his peanuts, Betty eats $9$ of her peanuts, and Charlie eats $25$ of his peanuts. Now the three numbers of peanuts each person has forms an arithmetic progression. Find the number of peanuts Alex had initially.


 
 \textbf{Answer: }108



\textbf{Solution 1}

Let $r$ be the common ratio, where $r>1$. We then have $ar-9-(a-5)=a(r-1)-4=ar^{2}-25-(ar-9)=ar(r-1)-16$. We now have, letting, subtracting the 2 equations, $ar^{2}+-2ar+a=12$, so we have $3ar=432,$ or $ar=144$, which is how much Betty had. Now we have $144+\dfrac{144}{r}+144r=444$, or $144(r+\dfrac{1}{r})=300$, or $r+\dfrac{1}{r}=\dfrac{25}{12}$, which solving for $r$ gives $r=\dfrac{4}{3}$, since $r>1$, so Alex had $\dfrac{3}{4} \cdot 144=\boxed{108}$ peanuts.

 





\textbf{Solution 2 (Quadratic Formula)}

Let $a$ be Alex's peanuts and $k$ the common ratio. Then we have $a(k^2+k+1)=444$. Adding $k$ to both sides and factoring,

 
 \[\frac{444}{a}+k=(k+1)^2\]

 For the common difference, $ak=5-(a-5)=ak^2-25-(ak-9)$. Simplifying, $k^2-2k+1=\frac{12}{a}$. Factoring, 
 \[(k-1)^2=\frac{12}{a}\]

 
 \[(k+1)^2-(k-1)^2=4k \implies 4k=\frac{444-12}{a}+k \implies k=\frac{144}{a}\]

 
Substitute $k$ in the second equation to get $(\frac{144-a}{a})^2=\frac{12}{a}$. Expanding and applying the quadratic formula, 
 \[a=150\pm\frac{\sqrt{300^2-4(144^2)}}{2}\]Taking out $4^2\cdot3^2$ from under the radical leaves 
 \[a=150\pm6\sqrt{625-576}=108, 192\]Since Alex's peanut number was the lowest of the trio, and $3*192>444$, Alex initially had $\boxed{108}$ peanuts.

 





\textbf{Solution 3}

Let the initial numbers of peanuts Alex, Betty and Charlie had be $a$, $b$, and $c$ respectively.Let the final numbers of peanuts, after eating, be $a'$, $b'$, and $c'$.

 We are given that $a + b + c = 444$. Since a total of $5 + 9 + 25 = 39$ peanuts are eaten, we must have $a' + b' + c' = 444 - 39 = 405$.Since $a'$, $b'$, and $c'$ form an arithmetic progression, we have that $a' = b' - x$ and $c' = b' + x$ for some integer $x$.Substituting yields $3b' = 405$ and so $b' = 135$. Since Betty ate $9$ peanuts, it follows that $b = b' + 9 = 144$.

 Since $a$, $b$, and $c$ form a geometric progression, we have that $a = \frac{b}{r}$ and $c = br$.Multiplying yields $ac = b^2 = 144^2$.Since $a + c = 444 - b = 300$, it follows that $a = 150 - \lambda$ and $c = 150 + \lambda$ for some integer $\lambda$.Substituting yields $(150-\lambda)(150+\lambda) = 144^2$, which expands and rearranges to $\lambda^2 = 150^2-144^2 = 42^2$.Since $\lambda > 0$, we must have $\lambda = 42$, and so $a = 150 - \lambda = \boxed{108}$.

 





\textbf{Solution 4}

Bashing is not difficult. All we have to consider is the first equation. We can write it as $x*(1+r+r^2) = 444$. The variable $x$ must be an integer, and after trying all the factors of $444$, it's clear that $r$ is a fraction smaller than $10$. When calculating the coefficient of $x$, we must consider that the fraction produced will very likely have a numerator that divides $444$. Trying a couple will make it easier to find the fraction, and soon you will find that $\frac{4}{3}$ gives a numerator of $37$, a rather specific factor of $444$. Solving for the rest will give you an integer value of $\boxed{108}$. This is by no means a good solution, but it may be faster in a competition if you don't want to mess with several other equations. This is purely up to different individuals.

 





\textbf{Solution 5}

Let $b$ be the finish number of Betty's peanuts. Then
 \[3b = 444-(5 + 9 + 25) = 405 = 3 \cdot 135 \implies b = 135, b+ 9 = 144.\]

 Let $k > 1$ be the common ratio. Then 
 \[\frac{144}{k} + 144 + k \cdot 144 = 444 \implies \frac{144}{k} + k \cdot 144 = 300\implies  \frac{12}{k} + k \cdot 12 = 25\implies k = \frac{4}{3} \implies \frac {144\cdot 3}{4} = \boxed {108}.\]

 \textbf{vladimir.shelomovskii@gmail.com, vvsss}

 




	\item (\textit{AIME II 2016}) There is a $40\%$ chance of rain on Saturday and a $30\%$ chance of rain on Sunday. However, it is twice as likely to rain on Sunday if it rains on Saturday than if it does not rain on Saturday. The probability that it rains at least one day this weekend is $\frac{a}{b}$, where $a$ and $b$ are relatively prime positive integers. Find $a+b$.


 
 \textbf{Answer: }107



\textbf{Solution 1}

Let $x$ be the probability that it rains on Sunday given that it doesn't rain on Saturday. We then have $\dfrac{3}{5}x+\dfrac{2}{5}2x = \dfrac{3}{10} \implies \dfrac{7}{5}x=\dfrac{3}{10}$ $\implies x=\dfrac{3}{14}$. Therefore, the probability that it doesn't rain on either day is $\left(1-\dfrac{3}{14}\right)\left(\dfrac{3}{5}\right)=\dfrac{33}{70}$. Therefore, the probability that rains on at least one of the days is $1-\dfrac{33}{70}=\dfrac{37}{70}$, so adding up the $2$ numbers, we have $37+70=\boxed{107}$.

 





\textbf{Solution 2}

Let Saturday $=6$ and Sunday $=7$. Then, $P(6) = 0.4$, $P(7) = 0.3$, and $P(7 \ | \ 6) = 2P(7 \ | \ 6^c)$. Consider the compliment, that is, $P(7^c \ | \ 6^c)$. This is exactly $1 - P(7 \ | \ 6^c)$. Note that $P(7) = P(7 \cap 6) + P(7 \cap 6^c)$ or $P(7) = P(7 \mid 6)P(6) + P(7 \mid 6^c)P(6^c)$. Let $P(7 \ | \ 6^c) = x$. Then, $P(7 \ | \ 6) = 2x$. So, $\frac{3}{10} = 2x(\frac{2}{5}) + x(\frac{3}{5}) \implies x = \frac{3}{14}$. We require $P(7^c \ | \ 6^c) = \frac{11}{14}$. Then, for $P(7^c \cup 6^c)$, we have $P(6^c) \cdot P(7^c \ | \ 6^c) = \frac{3}{5} \cdot \frac{11}{14} = \frac{33}{70}$. We need $P(7 \ | \ 6) \cup P(7 \ | \ 6^c) \cup P(7^c \ | \ 6) = 1 - P(7^c \cup 6^c) = 1 - \frac{33}{70} = \frac{37}{70} \implies \boxed{107}$.

 




	\item (\textit{AIME II 2016}) Let $x,y,$ and $z$ be real numbers satisfying the system

 \begin{align*} \log_2(xyz-3+\log_5 x)&=5,\\ \log_3(xyz-3+\log_5 y)&=4,\\ \log_4(xyz-3+\log_5 z)&=4. \end{align*}

Find the value of $|\log_5 x|+|\log_5 y|+|\log_5 z|$.


 
 \textbf{Answer: }265



\textbf{Solution}

First, we get rid of logs by taking powers: $xyz-3+\log_5 x=2^{5}=32$, $xyz-3+\log_5 y=3^{4}=81$, and $(xyz-3+\log_5 z)=4^{4}=256$. Adding all the equations up and using the $\log {xy}=\log {x}+\log{y}$ property, we have $3xyz+\log_5{xyz} = 378$, so we have $xyz=125$. Solving for $x,y,z$ by substituting $125$ for $xyz$ in each equation, we get $\log_5 x=-90, \log_5 y=-41, \log_5 z=134$, so adding all the absolute values we have $90+41+134=\boxed{265}$.

 Note: $xyz=125$ because we know $xyz$ has to be a power of $5$, and so it is not hard to test values in the equation $3xyz+\log_5{xyz} = 378$ in order to achieve desired value for $xyz$.

 




	\item (\textit{AIME II 2016}) An $a \times b \times c$ rectangular box is built from $a \cdot b \cdot c$ unit cubes. Each unit cube is colored red, green, or yellow. Each of the $a$ layers of size $1 \times b \times c$ parallel to the $(b \times c)$ faces of the box contains exactly $9$ red cubes, exactly $12$ green cubes, and some yellow cubes. Each of the $b$ layers of size $a \times 1 \times c$ parallel to the $(a \times c)$ faces of the box contains exactly $20$ green cubes, exactly $25$ yellow cubes, and some red cubes. Find the smallest possible volume of the box.


 
 \textbf{Answer: }180



\textbf{Solution}

By counting the number of green cubes $2$ different ways, we have $12a=20b$, or $a=\dfrac{5}{3} b$. Notice that there are only $3$ possible colors for unit cubes, so for each of the $1 \times b \times c$ layers, there are $bc-21$ yellow cubes, and similarly there are $ac-45$ red cubes in each of the $1 \times a \times c$ layers. Therefore, we have $a(bc-21)=25b$ and $b(ac-45)=9a$. We check a few small values of $a,b$ and solve for $c$, checking $(a,b)=(5,3)$ gives $c=12$ with a volume of $180$, $(a,b)=(10,6)$ gives $c=6$ with a volume of $360$, and $(a,b)=(15,9)$ gives $c=4$, with a volume of $540$. Any higher $(a,b)$ will $ab>180$, so therefore, the minimum volume is $\boxed{180}$.

 





\textbf{Solution 2}

The total number of green cubes is given by $12a=20b\Longrightarrow a=\frac{5}{3}b$. 

 Let $r$ be the number of red cubes on each one of the $b$ layers then the total number of red cubes is $9a=br$. Substitute $a=\frac{5}{3}b$ gives $r=15$.

 Repeating the procedure on the number of yellow cubes $y$ on each of the $a$ layers gives $y=15$.

 Therefore $bc=9+12+15=36$ and $ac=15+20+25=60$. Multiplying yields $abc^2=2160$.

 Since $abc^2$ is fixed, $abc$ is minimized when $c$ is maximized, which occurs when $a$, $b$ are minimized (since each of $ac$, $bc$ is fixed). Thus $(a,b,c)=(3,5,12)\Longrightarrow abc=\boxed{180}$

 




	\item (\textit{AIME II 2016}) Triangle $ABC_0$ has a right angle at $C_0$. Its side lengths are pairwise relatively prime positive integers, and its perimeter is $p$. Let $C_1$ be the foot of the altitude to $\overline{AB}$, and for $n \geq 2$, let $C_n$ be the foot of the altitude to $\overline{C_{n-2}B}$ in $\triangle C_{n-2}C_{n-1}B$. The sum $\sum_{n=2}^\infty C_{n-2}C_{n-1} = 6p$. Find $p$.


 
 \textbf{Answer: }182



\textbf{Solution 1}

--WARNING: Note this solution seems to have some serious errors --

 Do note that by counting the area in 2 ways, the first altitude is $x = \frac{ab}{c}$. By similar triangles, the common ratio is $\rho = \frac{a}{c}$ for each height, so by the geometric series formula, we have 
 \begin{align}     6p=\frac{x}{1-\rho} = \frac{ab}{c-a}. \end{align}Writing $p=a+b+c$ and clearing denominators, we get 
 \[13a=6p .\]Thus $p=13q$, $a=6q$, and $b+c=7q$, i.e. $c=7q-b$. Plugging these into $(1)$, we get $78q(q-b)=6bq$, i.e., $14b=13q$. Thus $q=14r$ and $p=182r$, $b=13r$, $a=84r$, $c=85r$. Taking $r=1$ (since $a,b,c$ are relatively prime) we get $p=\boxed{182}$.

 





\textbf{Solution 2}

Note that by counting the area in 2 ways, the first altitude is $\dfrac{ab}{c}$. By similar triangles, the common ratio is $\dfrac{a}{c}$ for reach height, so by the geometric series formula, we have $6p=\dfrac{\dfrac{ab}{c}}{1-\dfrac{a}{c}}$. Multiplying by the denominator and expanding, the equation becomes $\dfrac{ab}{c}=6a+6b+6c-\dfrac{6a^2}{c}-\dfrac{6ab}{c}-6a$. Cancelling $6a$ and multiplying by $c$ yields $ab=6bc+6c^2-6a^2-6ab$, so $7ab = 6bc+6b^2$ and $7a=6b+6c$. Checking for Pythagorean triples gives $13,84,$ and $85$, so $p=13+84+85=\boxed{182}$

 Note: a more rigorous solution instead of checking for triples would be to substitute $c = \sqrt{a^2 + b^2}$ and heavily simplifying the equation. Eventually we are left with $13a = 84b,$ and since $a, b$ are relatively prime, we know $a = 84$ and $b = 13.$ From here we can note that $(13, 84, 85)$ is a pythagorean triple or again use the Pythagorean Theorem to find $c = 85.$ Thus, the answer is $\boxed{182}.$

 Solution modified/fixed from Shaddoll's solution.

 





\textbf{Solution 3}

We start by splitting the sum of all $C_{n-2}C_{n-1}$ into two parts: those where $n-2$ is odd and those where $n-2$ is even.

 
First consider the sum of the lengths of the segments for which $n-2$ is odd for each $n\geq2$. The perimeters of these triangles can be expressed using $p$ and ratios that result because of similar triangles. Considering triangles where $n-2$ is odd, we find that the perimeter for each such $n$ is $p\left(\frac{C_{n-1}C_{n}}{C_{0}B}\right)$. Thus, 

 $p\sum_{n=1}^{\infty}\frac{C_{2n-1}C_{2n}}{C_{0}B}=6p+C_{0}B$.

 
Simplifying,

 $\sum_{n=1}^{\infty}C_{2n-1}C_{2n}=6C_{0}B + \frac{(C_{0}B)^2}{p}=C_{0}B\left(6+\frac{C_{0}B}{p}\right)$.     (1)

 
Continuing with a similar process for the sum of the lengths of the segments for which $n-2$ is even,

 $p\sum_{n=1}^{\infty}\frac{C_{2n-2}C_{2n-1}}{C_{0}B}=6p+C_{0}A+AB=7p-C_{0}B$.

 
Simplifying,

 $\sum_{n=1}^{\infty}C_{2n-2}C_{2n-1}=C_{0}B\left(7-\frac{C_{0}B}{p}\right)$.     (2)

 
Adding (1) and (2) together, we find that 

 $6p=13C_{0}B \Rightarrow p=\frac{13C_{0}B}{6}=C_{0}B+C_{0}A+AB \Rightarrow \frac{7C_{0}B}{6}=C_{0}A+AB \Rightarrow 7C_{0}B=6C_{0}A + 6AB$.

 Setting $a=C_{0}B$, $b=C_{0}A$, and $c=AB$, we can now proceed as in Shaddoll's solution, and our answer is $p=13+84+85=\boxed{182}$.

 





\textbf{Solution 4}


                    \begin{figure}[H]
                        \centering
                        \includegraphics[width=0.5\textwidth]{Solution/AIME_2016_II_5_Solution_Figure_1.png}
                    \end{figure}
                    

 Let $a = BC_0$, $b = AC_0$, and $c = AB$.Note that the total length of the red segments in the figure above is equal to the length of the blue segment times $\frac{a+c}{b}$.

 The desired sum is equal to the total length of the infinite path $C_0 C_1 C_2 C_3 \cdots$, shown in red in the figure below.Since each of the triangles $\triangle C_0 C_1 C_2, \triangle C_2 C_3 C_4, \dots$ on the left are similar, it follows that the total length of the red segments in the figure below is equal to the length of the blue segment times $\frac{a+c}{b}$.In other words, we have that $a\left(\frac{a+c}{b}\right) = 6p$.

 Guessing and checking Pythagorean triples reveals that $a = 84$, $b=13$, $c = 85$, and $p = a + b + c = \boxed{182}$ satisfies this equation.

 
                    \begin{figure}[H]
                        \centering
                        \includegraphics[width=0.5\textwidth]{Solution/AIME_2016_II_5_Solution_Figure_2.png}
                    \end{figure}
                    

 





\textbf{Solution 5}

This solution proceeds from $\frac{\frac{ab}{c}}{1-\frac{a}{c}} = \frac{\frac{ab}{c}}{\frac{c-a}{c}} = \frac{ab}{c-a} = 6(a+b+c)$. Note the general from for a primitive pythagorean triple, $m^2-n^2, 2mn, m^2+n^2$ and after substitution, letting $a = m^2-n^2, b = 2mn, c = m^2+n^2$ into the previous equation simplifies down very nicely into $m = 13n$. Thus $a = 168n^2, b = 26n^2, c = 170n^2$. Since we know all three side lengths are relatively prime, we must divide each by 2 and let n = 1 giving $a = 84, b = 13, c = 85$ yielding $p = a + b + c = \boxed{182}$.

 Alternate way to find a, b, and c (from VarunGotem) :  $7a=6b+6c$ implies that $6b=7a-6c$. Because $a^2+b^2=c^2$, $36a^2+(6b)^2=36c^2$. Plugging in gives $85a^2-84ac+36c^2=36c^2$. Simplifying gives $85a-84c=0$, and since $a$ and $c$ are relatively prime, $a=84$ and $c=85$. This means $b=13$ and $p=13+84+85=\boxed{182}$.

 





\textbf{Solution 6}

For this problem, first notice that its an infinite geometric series of $6(a+b+c)=\frac{ab}{c-b}$ if $c$ is the hypotenuse. WLOG $a<b$, we can generalize a pythagorean triple of $x^2-y^2, 2xy, x^2+y^2$. Let $b=2xy$, then this generalization gives $6(a+b+c)(c-b)=ab$
 \[(x^2-y^2)2xy=6(2x^2+2xy)(x-y)^2\]
 \[(x+y)xy=6(x^2+xy)(x-y)\]
 \[xy=6x(x-y)\]
 \[7xy=6x^2\]
 \[y=\frac{6}{7}x\]

 Now this is just clear. Let $x=7m$ and $y=6m$ for $m$ to be a positive integer, the pythagorean triple is $13-84-85$ which yields $\boxed{182}$.

 




	\item (\textit{AIME II 2016}) For polynomial $P(x)=1-\dfrac{1}{3}x+\dfrac{1}{6}x^{2}$, define
$Q(x)=P(x)P(x^{3})P(x^{5})P(x^{7})P(x^{9})=\sum_{i=0}^{50} a_ix^{i}$.
Then $\sum_{i=0}^{50} |a_i|=\dfrac{m}{n}$, where $m$ and $n$ are relatively prime positive integers. Find $m+n$.


 
 \textbf{Answer: }275



\textbf{Solution 1}

Note that $P(n)$ substitutes $n$ as $x$ into $P(x)$. So, $P(x^3)$ substitutes $x^3$, $P(x^5)$ substitutes $x^5$, etc. Notice how this doesn't change the coefficients of $P(x)$. So, the sum of the absolute values of the coefficients of $P(x)$, $P(x^3)$, $P(x^5)$, $P(x^7)$, $P(x^9)$ are all the same. The sum of the absolute values of the coefficients of $P(x)$ is $\frac{1}{6} + \frac{1}{3} + 1 = \frac{3}{2}$. Therefore, the sum of the coefficients of $Q(x)$ is simply $\left(\frac{3}{2}\right)\left(\frac{3}{2}\right)\left(\frac{3}{2}\right)\left(\frac{3}{2}\right)\left(\frac{3}{2}\right) = \left(\frac{3}{2}\right)^5 = \frac{243}{32}$, in which we have $243+32 = \boxed{275}$.

 





\textbf{Solution 2}

Note that all the coefficients of odd-powered terms is an odd number of odd degree terms multiplied together, and all coefficients of even-powered terms have an even number of odd degree terms multiplied together. Since every odd degree term is negative, and every even degree term is positive, the sum is just equal to $Q(-1)=P(-1)^{5}=\left( \dfrac{3}{2}\right)^{5}=\dfrac{243}{32}$, so the desired answer is $243+32=\boxed{275}$.

 [YINGLINWU] Just a quick comment: After multiplying out the product for Q(x) into the sum of monomials, the coefficient of a monomial of even degree is always positive while the coefficient of a monomial of odd degree is always negative. Therefore, as we are collecting similar terms, no coefficients of monomials cancel out. It follows that the sum of absolute values of coefficients of Q(x) is the same as that of Q(-x), which is given by Q(-1).

 





\textbf{Solution 3}

We are looking for the sum of the absolute values of the coefficients of $Q(x)$. By defining $P'(x) = 1 + \frac{1}{3}x+\frac{1}{6}x^2$, and defining $Q'(x) = P'(x)P'(x^3)P'(x^5)P'(x^7)P'(x^9)$, we have made it so that all coefficients in $Q'(x)$ are just the positive/absolute values of the coefficients of $Q(x)$. . 

 
To find the sum of the absolute values of the coefficients of $Q(x)$, we can just take the sum of the coefficients of $Q'(x)$. This sum is equal to
 \[Q'(1) = P'(1)P'(1)P'(1)P'(1)P'(1) = \left(1+\frac{1}{3}+\frac{1}{6}\right)^5 = \frac{243}{32},\]

 so our answer is $243+32 = \boxed{275}$.

 Note: this method doesn't work for every product of polynomials. Example: The sum of the absolute values of the coefficients of $K(x) = (2x^2 - 6x - 5)(10x - 1)$ is $131$ but when you find the sum of the coefficients of $K'(x)$ which is $K'(1)$, then you get $143$. - whatRthose

 





\textbf{Solution 4 (risky)}

Multiply $P(x)P(x^3)$ and notice that the odd degree terms have a negative coefficient. Observing that this is probably true for all polynomials like this (including $P(x)P(x^3)P(x^5)P(x^7)P(x^9)$), we plug in $-1$ to get $\frac{243}{32} \implies \boxed{275}$.

 





\textbf{Solution 5 (bash for life)}

We expand and add the numbers (I made two very easy mistakes and fixed them).

 \href{https://cdn.artofproblemsolving.com/attachments/1/6/c4e3bea5cf2d5bb7d7796049b5dc17fcd8a2ba.jpg}{https://cdn.artofproblemsolving.com/attachments/1/6/c4e3bea5cf2d5bb7d7796049b5dc17fcd8a2ba.jpg}

 -maxamc

 




	\item (\textit{AIME II 2016}) Squares $ABCD$ and $EFGH$ have a common center and $\overline{AB} || \overline{EF}$. The area of $ABCD$ is 2016, and the area of $EFGH$ is a smaller positive integer. Square $IJKL$ is constructed so that each of its vertices lies on a side of $ABCD$ and each vertex of $EFGH$ lies on a side of $IJKL$. Find the difference between the largest and smallest positive integer values for the area of $IJKL$.


 
 \textbf{Answer: }840



\textbf{Solution}

Letting $AI=a$ and $IB=b$, we have 
 \[IJ^{2}=a^{2}+b^{2} \geq 1008\] by \href{https://artofproblemsolving.com/wiki/index.php?title=AM-GM\_inequality}{AM-GM inequality}. Also, since $EFGH||ABCD$, the angles that each square cuts another are equal, so all the triangles are formed by a vertex of a larger square and $2$ adjacent vertices of a smaller square are similar. Therefore, the areas form a geometric progression, so since 
 \[2016=12^{2} \cdot 14\] we have the maximum area is 
 \[2016 \cdot \dfrac{11}{12} = 1848\] (the areas of the squares from largest to smallest are $12^{2} \cdot 14, 11 \cdot 12 \cdot 14, 11^{2} \cdot 14$ forming a geometric progression). 

 
The minimum area is $1008$ (every square is half the area of the square whose sides its vertices touch), so the desired answer is 
 \[1848-1008=\boxed{840}\]

 
                    \begin{figure}[H]
                        \centering
                        \includegraphics[width=0.5\textwidth]{Solution/AIME_2016_II_7_Solution_Figure_1.png}
                    \end{figure}
                    

 




	\item (\textit{AIME II 2016}) Find the number of sets $\{a,b,c\}$ of three distinct positive integers with the property that the product of $a,b,$ and $c$ is equal to the product of $11,21,31,41,51,$ and $61$.


 
 \textbf{Answer: }728



\textbf{Solution 1}

Note that the prime factorization of the product is $3^{2}\cdot 7 \cdot 11 \cdot 17 \cdot 31 \cdot 41 \cdot 61$. Ignoring overcounting, by stars and bars there are [2 (\# of 3's) plus 3 (a, b, c)  - 1] choose [3 - 1] = $6$ ways to choose how to distribute the factors of $3$, and $3$ ways to distribute the factors of the other primes, so we have $3^{6} \cdot 6$ ways. However, some sets have $2$ numbers that are the same, namely the ones in the form $1,1,x$ and $3,3,x$, which are each counted $3$ times, and each other set is counted $6$ times, so the desired answer is $\dfrac{729 \cdot 6-6}{6} = \boxed{728}$.

 





\textbf{Solution 2}

Again, notice that the prime factors of the product are $3, 3, 7, 11, 17, 31, 41, 61$. In this problem, we are asked to partition this set of distinct(ish) factors into three smaller indistinct sets. To do this, we can use Stirling numbers of the second kind, but Stirling numbers of the second kind would assume no empty parts, which isn't what we want. However, this is easy to fix. Denote Stirling numbers of the second kind with $S(n,k)$. We may start at the situation when all three sets are nonempty. Then, the number of partitions is simply $S(8,3)$. However, we are overcounting, since every time the two threes are in different sets, (unless they are both individually in their own sets), each one is double counted. To fix this, we can see that they were in different sets $S(8, 3) - S(7,3)$ times by complementary counting, but one of these times they were in their own individual sets, so the total overcount is $\dfrac{S(8,3) - S(7,3) - 1}{2}$. Similarly, we can do the cases for if two of the sets are nonempty and one of the sets are nonempty, (but in those last two cases the threes cannot be individually in their own sets). We then find that the "answer" is given by:
 \[S(8,3) - \dfrac{S(8,3) - S(7,3) - 1}{2} + S(8,2) - \dfrac{S(8,2) - S(7,2)}{2}+ S(8,1) - \dfrac{S(8,1) - S(7,1)}{2}\]
 \[= \dfrac{S(8,3)+S(8,2) + S(8,1) + S(7,3) + S(7,2) + S(7,1) + 1}{2}\]Drawing out the Stirling number triangle and evaluating yields $730$ for the last expression. However, throughout the entire solution, we ignored the fact that the three numbers $a$, $b$, and $c$ needed to be distinct. However, this is easy to fix, since there are only two sets in which they are not, namely $(1,1,x)$ and $(3,3,x)$. Thus, the actual answer is $730 - 2 = \boxed{728}$.

 - Aathreyakadambi

 




	\item (\textit{AIME II 2016}) The sequences of positive integers $1,a_2, a_3,...$ and $1,b_2, b_3,...$ are an increasing arithmetic sequence and an increasing geometric sequence, respectively. Let $c_n=a_n+b_n$. There is an integer $k$ such that $c_{k-1}=100$ and $c_{k+1}=1000$. Find $c_k$.


 
 \textbf{Answer: }262



\textbf{Solution 1}

Since all the terms of the sequences are integers, and 100 isn't very big, we should just try out the possibilities for $b_2$. When we get to $b_2=9$ and $a_2=91$, we have $a_4=271$ and $b_4=729$, which works, therefore, the answer is $b_3+a_3=81+181=\boxed{262}$.

 





\textbf{Solution 2 (Some trial and error)}

We have $a_k=r^{k-1}$ and $b_k=(k-1)d$. First, $b_{k-1}<c_{k-1}=100$ implies $d<100$, so $b_{k+1}<300$. 

 It follows that $a_{k+1}=1000-b_{k+1}>700$, i.e., 
 \[700 < r^k < 1000.\] Moreover, since $k$ is atleast $3$ we get $r^3\le r^k <1000$, i.e. $r<10$. For every value of $r$ in this range, define $i(r) = \max \{x : r^x < 700\}$, and define $j(r) = \min \{x : r^x > 1000\}$. We are looking for values of $r$ such that $j(r) -i(r)>1$. Let's make a table:
 \begin{array}[b]{ c c c c c c c c c } r & 2 & 3 & 4 & 5 & 6 & 7 & 8 & 9 \\[2ex] i(r) & 9  & 5 & 4  & 4 & 3 & 3 & 3 & 2\\[2ex] j(r) & 10  & 7 & 5  & 5 & 4 & 4 & 4 & 4 \end{array}The only admissible values for $r^k$ are $\{3^6, 9^3\}$. However, since $100=c_{k-1}=r^{k-2}+(k-2)d+1$, we must have $(k-2)\mid 99-r^{k-2}$. This does not hold for $r^k=3^6$ because $4$ does not divide $99-3^4=18$. This leaves $r^k=9^3$ as the only option. 

 For $r=9$ and $k=3$, we check: $a_{k-1}= a_2 = r =9$ implies $b_{k-1}= b_2 = 91$, i.e. $d=90$. Then $a_{k+1}=a_4 = r^3 = 729$ and $b_{k+1}=b_4 = 1+3d = 271$ and $c_{k+1}=c_4=a_4+b_4 = 729+271=1000$; so it works! Then $c_k = c_3 = 9^2+181 = 262$.

 





\textbf{Solution 3}

Using the same reasoning ($100$ isn't very big), we can guess which terms will work. The first case is $k=3$, so we assume the second and fourth terms of $c$ are $100$ and $1000$. We let $r$ be the common ratio of the geometric sequence and write the arithmetic relationships in terms of $r$.

 The common difference is $100-r - 1$, and so we can equate: $2(99-r)+100-r=1000-r^3$. Moving all the terms to one side and the constants to the other yields

 $r^3-3r = 702$, or $r(r^2-3) = 702$. Simply listing out the factors of $702$ shows that the only factor $3$ less than a square that works is $78$. Thus $r=9$ and we solve from there to get $\boxed{262}$.

 Solution by rocketscience

 





\textbf{Solution 4 (More Robust Bash)}

The reason for bashing in this context can also be justified by the fact 100 isn't very big. 

 Let the common difference for the arithmetic sequence be $a$, and the common ratio for the geometric sequence be $b$. The sequences are now $1, a+1, 2a+1, \ldots$, and $1, b, b^2, \ldots$. We can now write the given two equations as the following:

 $1+(k-2)a+b^{k-2} = 100$

 $1+ka+b^k = 1000$

 Take the difference between the two equations to get $2a+(b^2-1)b^{k-2} = 900$. Since 900 is divisible by 4, we can tell $a$ is even and $b$ is odd. Let $a=2m$, $b=2n+1$, where $m$ and $n$ are positive integers. Substitute variables and divide by 4 to get:

 $m+(n+1)(n)(2n+1)^{k-2} = 225$

 Because very small integers for $n$ yield very big results, we can bash through all cases of $n$. Here, we set an upper bound for $n$ by setting $k$ as 3. After trying values, we find that $n\leq 4$, so $b=9, 7, 5, 3$. Testing out $b=9$ yields the correct answer of $\boxed{262}$. Note that even if this answer were associated with another b value like $b=3$, the value of $k$ can still only be 3 for all of the cases.

 





\textbf{Solution 5 (Casework)}

Let $a_n$ and $b_n$ be in the form of
 \begin{array}[b]{ c c c c c c c } n & 1 & 2 & 3 & 4 & 5 & 6 \\[2ex] a_n & 1 & a+1 & 2a+1 & 3a+1 & 4a+1 & 5a+1 \\[2ex] b_n & 1 & b & b^2 & b^3 & b^4 & b^5 \\[2ex] c_n & 2 & b+a+1 & b^2+2a+1 & b^3+3a+1 & b^4+4a+1 & b^5+5a+1 \end{array}Case $1.\hspace{10mm}   k = 3\hspace{5mm} (c_1=2 \implies k > 2).$
 \[c_2 = a+1 + b = 100, c_4 = 3a+1 + b^3 = 1000 \implies b^3 -3b -2 = 1000-300 \implies   b^3 - 3b = 702 \implies b = 9, a=90, c_3 = \boxed {262}.\]Case $2. \hspace{10mm}   k = 4.$
 \[c_3 = 2a+1 + b^2 = 100, c_5 = 4a+1 + b^4 = 1000 \implies b^4 -2b^2 -1 = 1000-200 \implies   b^4 - 2b^2 = 801 \implies \O.\]Case $3. \hspace{10mm}   k \ge 5.\hspace{3mm}  b^5 < 1000 \implies b = \{2,3\}.$

 Case $3.1.\hspace{10mm} b = 2.$
 \[c_{k-1} = 2^{k-2} + (k-2) a + 1 = 100, c_{k+1} = 2^k + ka + 1 = 1000\implies a = 450-3\cdot 2^{k-3} \implies 2^k +450k -3k\cdot 2^{k-3} + 1 = 1000  \implies \O.\]Case $3.2.\hspace{10mm} b = 3.$
 \[c_{k-1} = 3^{k-2} + (k-2) a + 1 = 100, c_{k+1} = 3^k + ka + 1 = 1000\implies a = 450-4\cdot 3^{k-2} \implies 3^{k-4}(9-4k) +50k = 3\cdot 37  \implies \O.\]\textbf{vladimir.shelomovskii@gmail.com, vvsss}

 




	\item (\textit{AIME II 2016}) Triangle $ABC$ is inscribed in circle $\omega$. Points $P$ and $Q$ are on side $\overline{AB}$ with $AP<AQ$. Rays $CP$ and $CQ$ meet $\omega$ again at $S$ and $T$ (other than $C$), respectively. If $AP=4,PQ=3,QB=6,BT=5,$ and $AS=7$, then $ST=\frac{m}{n}$, where $m$ and $n$ are relatively prime positive integers. Find $m+n$.


 
 \textbf{Answer: }043



\textbf{Solution 1}


                    \begin{figure}[H]
                        \centering
                        \includegraphics[width=0.5\textwidth]{Solution/AIME_2016_II_10_Solution_Figure_1.png}
                    \end{figure}
                    Let $\angle ACP=\alpha$, $\angle PCQ=\beta$, and $\angle QCB=\gamma$. Note that since $\triangle ACQ\sim\triangle TBQ$ we have $\frac{AC}{CQ}=\frac56$, so by the Ratio Lemma 
 \[\dfrac{AP}{PQ}=\dfrac{AC}{CQ}\cdot\dfrac{\sin\alpha}{\sin\beta}\quad\implies\quad \dfrac{\sin\alpha}{\sin\beta}=\dfrac{24}{15}.\]Similarly, we can deduce $\frac{PC}{CB}=\frac47$ and hence $\frac{\sin\beta}{\sin\gamma}=\frac{21}{24}$.

 Now Law of Sines on $\triangle ACS$, $\triangle SCT$, and $\triangle TCB$ yields 
 \[\dfrac{AS}{\sin\alpha}=\dfrac{ST}{\sin\beta}=\dfrac{TB}{\sin\gamma}.\]Hence 
 \[\dfrac{ST^2}{\sin^2\beta}=\dfrac{TB\cdot AS}{\sin\alpha\sin\gamma},\]so 
 \[TS^2=TB\cdot AS\left(\dfrac{\sin\beta}{\sin\alpha}\dfrac{\sin\beta}{\sin\gamma}\right)=\dfrac{15\cdot 21}{24^2}\cdot 5\cdot 7=\dfrac{35^2}{8^2}.\]Hence $ST=\frac{35}8$ and the requested answer is $35+8=\boxed{43}$.

 Edit: Note that the finish is much simpler. Once you get $\dfrac{AS}{\sin\alpha}=\dfrac{ST}{\sin\beta}$, you can solve quickly from there getting $ST=\dfrac{AS \sin(\beta)}{\sin(\alpha)}=7\cdot \dfrac{15}{24}=\dfrac{35}{8}$.

 





\textbf{Solution 2 (Projective Geometry)}


                    \begin{figure}[H]
                        \centering
                        \includegraphics[width=0.5\textwidth]{Solution/AIME_2016_II_10_Solution_Figure_2.png}
                    \end{figure}
                    Projecting through $C$ we have 
 \[\frac{3}{4}\times \frac{13}{6}=(A,Q;P,B)\stackrel{C}{=}(A,T;S,B)=\frac{ST}{7}\times \frac{13}{5}\] which easily gives $ST=\frac{35}{8}\Longrightarrow 35+8=\boxed{043}$.

 





\textbf{Solution 3}

By Ptolemy's Theorem applied to quadrilateral $ASTB$, we find
 \[5\cdot 7+13\cdot ST=AT\cdot BS.\]Therefore, in order to find $ST$, it suffices to find $AT\cdot BS$. We do this using similar triangles, which can be found by using Power of a Point theorem.

 As $\triangle APS\sim \triangle CPB$, we find
 \[\frac{4}{PC}=\frac{7}{BC}.\]Therefore, $\frac{BC}{PC}=\frac{7}{4}$.

 As $\triangle BQT\sim\triangle CQA$, we find
 \[\frac{6}{CQ}=\frac{5}{AC}.\]Therefore, $\frac{AC}{CQ}=\frac{5}{6}$.

 As $\triangle ATQ\sim\triangle CBQ$, we find
 \[\frac{AT}{BC}=\frac{7}{CQ}.\]Therefore, $AT=\frac{7\cdot BC}{CQ}$.

 As $\triangle BPS\sim \triangle CPA$, we find
 \[\frac{9}{PC}=\frac{BS}{AC}.\]Therefore, $BS=\frac{9\cdot AC}{PC}$. Thus we find
 \[AT\cdot BS=\left(\frac{7\cdot BC}{CQ}\right)\left(\frac{9\cdot AC}{PC}\right).\]But now we can substitute in our previously found values for $\frac{BC}{PC}$ and $\frac{AC}{CQ}$, finding
 \[AT\cdot BS=63\cdot \frac{7}{4}\cdot \frac{5}{6}=\frac{21\cdot 35}{8}.\]Substituting this into our original expression from Ptolemy's Theorem, we find
 \begin{align*}35+13ST&=\frac{21\cdot 35}{8}\\13ST&=\frac{13\cdot 35}{8}\\ST&=\frac{35}{8}.\end{align*}
Thus the answer is $\boxed{43}$.

 





\textbf{Solution 4}

Extend $\overline{AB}$ past $B$ to point $X$ so that $CPTX$ is cyclic. Then, by Power of a Point on $CPTX$, $(CQ)(QT) =  (PQ)(QX)$. By Power of a Point on $CATB$, $(CQ)(QT) = (AQ)(QB) = 42$. Thus, $(PQ)(QX) = 42$, so $BX = 8$. 

 By the Inscribed Angle Theorem on $CPTX$, $\angle SCT = \angle BXT$. By the Inscribed Angle Theorem on $ASTC$, $\angle SCT = \angle SAT$, so $\angle BXT = \angle SAT$. Since $ASTB$ is cyclic, $\angle AST = \angle TBX$. Thus, $\triangle AST \sim \triangle XBT$, so $AS/XB = ST/BT$. Solving for $ST$ yields $ST = \frac{35}{8}$, for a final answer of $35+8 = \boxed{043}$.

 





\textbf{Solution 5 (5 = 2 + 3)}


                    \begin{figure}[H]
                        \centering
                        \includegraphics[width=0.5\textwidth]{Solution/AIME_2016_II_10_Solution_Figure_3.png}
                    \end{figure}
                    By Ptolemy's Theorem applied to quadrilateral $ASTB$, we find
 \[AS\cdot BT+AB\cdot ST=AT\cdot BS.\]Projecting through $C$ we have
 \[\frac{AQ \cdot PB}{PQ \cdot AB} = (A,Q; P,B)\stackrel{C}{=}(A,T; S,B)=\frac{AT \cdot BS}{ST \cdot AB}.\]Therefore  
 \[AT \cdot BS = \frac {AQ \cdot PB}{PQ} \times ST \implies\]
 \[\left(\frac {AQ \cdot PB}{PQ} - AB\right)\times ST = AS \cdot BT \implies\]
 \[ST = \frac {AS \cdot BT \cdot PQ}{AQ \cdot PB – AB \cdot PQ}\]
 \[ST = \frac {7\cdot 5 \cdot 3}{7\cdot 9 – 13 \cdot 3 } = \frac {35}{8} \implies 35 + 8 = \boxed {43}.\]

 \textbf{vladimir.shelomovskii@gmail.com, vvsss}

 





\textbf{Solution 6}

Connect $AT$ and $\angle{SCT}=\angle{SAT}, \angle{ACS}=\angle{ATS}, \frac{ST}{\sin \angle{SAT}}=\frac{AS}{\sin \angle{ATS}}$

 So we need to get the ratio of $\frac{\sin \angle{ACS}}{\sin \angle{SCT}}$

 By clear observation $\triangle{CAQ}\sim \triangle{BTQ}$, we have $\frac{CQ}{AC}=\frac{6}{5}$, LOS tells $\frac{AC}{\sin \angle{CPA}}=\frac{4}{\sin \angle{ACS}}; \frac{CQ}{\sin \angle{CPQ}}=\frac{3}{\sin \angle{PCQ}}$ so we get $\frac{\sin \angle{PCQ}}{\sin \angle{ACS}}=\frac{5}{8}$, the desired answer is $7\cdot \frac{\sin \angle{SAT}}{\sin \angle{ATS}}=\frac{35}{8}$ leads to $\boxed{043}$

 





\textbf{Solution 7 (no trig or projections)}

Note that since $\triangle SAP~\triangle BCP$, $\frac{9}{SP}=\frac{BC}{7}=\frac{PC}{4}$. Furthermore, since $\triangle ACQ~\triangle TBQ$, we have $\frac{7}{TQ}=\frac{AC}{5}=\frac{QC}{6}$. From Stewart's on triangle $BCP$, we have $25CQ+BC^2\cdot TQ=TQ\cdot CQ\cdot TC+36TC$, and since $TQ\cdot CQ=6\cdot7=42$ by power of a point, this simplifies to $25CQ+BC^2\cdot TQ=78TC$. Similarly, $49CP+AC^2\cdot SP=52SC$. Finally, using Ptolemy's on quadrilateral $ACBS$ yields $13SC=7BC+SB\cdot AC$, and using Ptolemy's on quadrilateral $ACBT$ yields $13TC=5AC+TA\cdot BC$. From Ptolemy's on $ABTS$, we find $SB\cdot TA=13ST+35$, which is nice because it contains $ST$.We return to our first Stewart's equation: $25CQ+BC^2\cdot TQ=78TC$, and we notice that $CQ$ and $TQ$ can be related to $AC$ using our similar triangle conditions. Substituting gives us $30AC+\frac{35BC^2}{AC}=78TC$, which by four times our first Ptolemy's equation also equals $30AC+6TA\cdot BC$. Thus, $\frac{35BC^2}{AC}=6TA\cdot BC$ and $TA=\frac{35}{6}\cdot\frac{BC}{AC}$. Similarly, from our other Stewart's equation, we find $28BC+\frac{63AC^2}{BC}=52SC=28BC+4SB\cdot AC$, or $SB=\frac{63}{4}\cdot\frac{AC}{BC}$. Plugging this into our final Ptolemy's equation, we find 
 \[SB\cdot TA=13ST+35\Longrightarrow\frac{35\cdot63}{6\cdot4}=13ST+35\Longrightarrow ST=\frac{\frac{35\cdot21}{8}-35}{13}=\frac{35\cdot\frac{13}{8}}{13}=\frac{35}{8},\]giving us our final answer of $\boxed{043}$.

 





\textbf{Solution 8(Algebra Bash)}

Let $\angle STB = \theta$. Since quadrilateral $STBA$ is cyclic, let $ST = x$ and we get

 $x^2 + 25 - 10x \cos(\theta) = 49 + 169 - 182 \cos(\theta) \implies \cos(\theta) = \frac{x^2 - 193}{10x + 182}$

 so if we figure out $\cos(\theta)$ we're done as we can easily solve for $x$. 

 Now let $PC = a$ and $QC = b$. Note $\triangle SPA \sim \triangle QPC \implies BC = \frac{7a}{4}$

 Also $\triangle BQT \sim \triangle AQC \implies AC = \frac{5b}{6}$. 

 Now we use Stewart's Theorem on $\triangle AQC, \triangle BPC$. We respectively get

 $4 \cdot 7 \cdot 3 + a^2 \cdot 7 = \frac{25b^2}{36} \cdot 3 + b^2 \cdot 4 \implies 84 + 7a^2 = b^2(\frac{73}{12})$

 $3 \cdot 9 \cdot 6 + 9 \cdot b^2 = 6 \cdot a^2 + 3 \cdot \frac{49a^2}{16} \implies b^2 = \frac{27a^2}{16} - 18$

 Thus, we solve

 $84 + 7a^2 = (\frac{27a^2}{16} - 18)(\frac{73}{12})$ so we get $a^2 = \frac{12384}{209}$

 Then we use $\frac{65a^2}{16} + \frac{7a^2}{2} \cos(\theta) = 81$ to solve for $\cos(\theta)$ and we get

 $\cos(\theta) = \frac{-3709}{4816}$

 Finally, we solve $\frac{x^2 - 193}{10x + 182} = \cos(\theta) = \frac{-3709}{4816} \implies x = \frac{35}{8}, \frac{-3635}{301}$ and obviously the answer is $\frac{35}{8} \implies \boxed{043}$

 





\textbf{Solution 9 (Another Projective)}

All lengths are directed. We have 
 \[(B, S; T, A) \stackrel{C}{=} (B,P; Q,A) = -\frac{8}{13}\]

 Therefore, 
 \[\frac{TB}{TS} \div \frac{AB}{AS} = -\frac{8}{13}\]

 Solving yields that $ST = \frac{8}{35}$, so our answer is $\boxed {043}$.

 




	\item (\textit{AIME II 2016}) For positive integers $N$ and $k$, define $N$ to be $k$-nice if there exists a positive integer $a$ such that $a^{k}$ has exactly $N$ positive divisors. Find the number of positive integers less than $1000$ that are neither $7$-nice nor $8$-nice.


 
 \textbf{Answer: }749



\textbf{Solution}

We claim that an integer $N$ is only $k$-nice if and only if $N \equiv 1 \pmod k$. By the number of divisors formula, the number of divisors of $\prod_{i=1}^n p_i^{a_i}$ is $\prod_{i=1}^n (a_i+1)$. Since all the $a_i$'s are divisible by $k$ in a perfect $k$ power, the only if part of the claim follows. To show that all numbers $N \equiv 1 \pmod k$ are $k$-nice, write $N=bk+1$. Note that $2^{kb}$ has the desired number of factors and is a perfect kth power. By PIE, the number of positive integers less than $1000$ that are either $1 \pmod 7$ or $1\pmod 8$ is $143+125-18=250$, so the desired answer is $999-250=\boxed{749}$.

 Solution by Shaddoll and firebolt360

 





\textbf{Solution 2}

All integers $a$ will have factorization $2^a3^b5^c7^d...$. Therefore, the number of factors in $a^7$ is $(7a+1)(7b+1)...$, and for $a^8$ is $(8a+1)(8b+1)...$. The most salient step afterwards is to realize that all numbers $N$ not $1 \pmod{7}$ and also not $1 \pmod{8}$ satisfy the criterion. The cycle repeats every $56$ integers, and by PIE, $7+8-1=14$ of them are either $7$-nice or $8$-nice or both. Therefore, we can take $\frac{42}{56} * 1008 = 756$ numbers minus the $7$ that work between $1000-1008$ inclusive, to get $\boxed{749}$ positive integers less than $1000$ that are not nice for $k=7, 8$.

 





\textbf{Solution 3}

An easier way to get $N \equiv 1 \pmod k$. Notice that they define $k$-nice to be $a^k$ has $N$ positive divisors. This implies that if $a$ is prime, then $N$ is only $k$-nice if and only if $k+1 = N$. Thus, we have proved that $N \equiv 1 \pmod k$ for $a$ prime. Then, assume $a$ is composite. Then, $a$ itself can be split into multiple factors, each containing $kx+1$ per term for the prime factorization. Then, notice that this case becomes trivial! All factors split into multiples of $k$, and obviously sums of $k$ are divisible by $k$, EXCEPT for one, that being, 1 itself. Then, we have proved that $N$ is also $k$-nice if $N \equiv 1\pmod k$. Then, as said in prior solutions, we must not have $n \equiv 1 \pmod 7$ and $n \equiv 1 \pmod 8$. Counting all $1 \pmod 7$ between 1 and 1000, we have every number in the form $7N + 1$. There are then 143 numbers that satisfy that constraint. Then, for $1 \pmod 8$, we have $8N + 1$, and there are 125 number that satisfy that. Then, to deal with $1 \pmod {56}$ (we over counted), we have 18 numbers, and thus to total number of $7$ and $8$ nice numbers is $143+125 - 18 = 250$ numbers, and our answer is $999-250 = \boxed{749}$.

 




	\item (\textit{AIME II 2016}) The figure below shows a ring made of six small sections which you are to paint on a wall. You have four paint colors available and you will paint each of the six sections a solid color. Find the number of ways you can choose to paint the sections if no two adjacent sections can be painted with the same color.


 
                    \begin{figure}[H]
                        \centering
                        \includegraphics[width=0.5\textwidth]{Problem/AIME_2016_II_12_Problem_Figure_1.png}
                    \end{figure}
                    


 
 \textbf{Answer: }732



\textbf{Solution 1}

Choose a section to start coloring. Assume, WLOG, that this section is color $1$. We proceed coloring clockwise around the ring. Let $f(n,C)$ be the number of ways to color the first $n$ sections (proceeding clockwise) such that the last section has color $C$. In general (except for when we complete the coloring), we see that
 \[f(n,C_i)=\sum_{j\ne i} f(n-1,C_j),\]i.e., $f(n,C_i)$ is equal to the number of colorings of $n-1$ sections that end in any color other than $C_i$. Using this, we can compute the values of $f(n,C)$ in the following table.

 $\begin{tabular}{c|c|c|c|c } \multicolumn{1}{c}{}&\multicolumn{4}{c}{\(C\)}\\ \(n\)&1 & 2 & 3& 4 \\ \hline 1& 1 & 0 & 0 & 0\\ 2 & 0 & 1 & 1 & 1 \\ 3& 3 & 2 & 2 & 2 \\ 4 & 6 & 7 & 7 & 7 \\ 5 & 21 & 20 & 20 & 20\\ 6& 0& 61 & 61 & 61\\ \end{tabular}$

 Note that $f(6, 1)=0$ because then $2$ adjacent sections are both color $1$. We multiply this by $4$ to account for the fact that the initial section can be any color.  Thus the desired answer is $(61+61+61) \cdot 4 = \boxed{732}$.

 Solution by Shaddoll

 





\textbf{Solution 2}

We use complementary counting. There are $4^6$ total colorings of the ring without restriction. To count the complement, we wish to count the number of colorings in which at least one set of adjacent sections are the same color. There are six possible sets of adjacent sections that could be the same color (think of these as borders). Call these $B_1,B_2,\dots,B_6$. Let $\mathcal{A}_1, \mathcal{A}_2,\dots,\mathcal{A}_6$ be the sets of colorings of the ring where the sections on both sides of $B_1,B_2,\dots,B_6$ are the same color. We wish to determine $|\mathcal{A}_1\cup\mathcal{A}_2\cup\cdots\cup\mathcal{A}_6|$. Note that all of these cases are symmetric, and in general, $|\mathcal{A}_i|=4^5$. There are $6$ such sets $\mathcal{A}_i$. Also, $|\mathcal{A}_i\cup\mathcal{A}_j|=4^4$, because we can only change colors at borders, so if we have two borders along which we cannot change colors, then there are four borders along which we have a choice of color. There are $\binom{6}{2}$ such pairs $\mathcal{A}_i\cup\mathcal{A}_j$. Similarly, $|\mathcal{A}_i\cup \mathcal{A}_j\cup \mathcal{A}_k|=4^3$, with $\binom{6}{3}$ such triples, and we see that the pattern will continue for 4-tuples and 5-tuples. For 6-tuples, however, these cases occur when there are no changes of color along the borders, i.e., each section has the same color. Clearly, there are four such possibilities.

 Therefore, by PIE,
 \[|\mathcal{A}_1\cup\mathcal{A}_2\cup\cdots\cup\mathcal{A}_6|=\binom{6}{1}\cdot 4^5-\binom{6}{2}\cdot 4^4+\binom{6}{3}\cdot 4^3-\binom{6}{4}\cdot 4^2+\binom{6}{5}\cdot 4^1-4.\]We wish to find the complement of this, or 
 \[4^6-\left(\binom{6}{1}\cdot 4^5-\binom{6}{2}\cdot 4^4+\binom{6}{3}\cdot 4^3-\binom{6}{4}\cdot 4^2+\binom{6}{5}\cdot 4^1-4\right).\]By the Binomial Theorem, this is $(4-1)^6+3=\boxed{732}$.

 





\textbf{Solution 3}

We use generating functions. Suppose that the colors are $0,1,2,3$.  Then as we proceed around a valid coloring of the ring in the clockwise direction, we know that between two adjacent sections with colors $s_i$ and $s_{i+1}$, there exists a number $d_i\in\{1,2,3\}$ such that $s_{i+1}\equiv s_i+d_i\pmod{4}$. Therefore, we can represent each border between sections by the generating function $(x+x^2+x^3)$, where $x,x^2,x^3$ correspond to increasing the color number by $1,2,3\pmod4$, respectively. Thus the generating function that represents going through all six borders is $A(x)=(x+x^2+x^3)^6$, where the coefficient of $x^n$ represents the total number of colorings where the color numbers are increased by a total of $n$ as we proceed around the ring. But if we go through all six borders, we must return to the original section, which is already colored. Therefore, we wish to find the sum of the coefficients of $x^n$ in $A(x)$ with $n\equiv 0\pmod4$.

 Now we note that if $P(x)=x^n$, then
 \[P(x)+P(ix)+P(-x)+P(-ix)=\begin{cases}4x^n&\text{if } n\equiv0\pmod{4}\\0&\text{otherwise}.\end{cases}\]Therefore, the sum of the coefficients of $A(x)$ with powers congruent to $0\pmod 4$ is
 \[\frac{A(1)+A(i)+A(-1)+A(-i)}{4}=\frac{3^6+(-1)^6+(-1)^6+(-1)^6}{4}=\frac{732}{4}.\]We multiply this by $4$ to account for the initial choice of color, so our answer is $\boxed{732}$.

 





\textbf{Solution 4}

Let $f(n)$ be the number of valid ways to color a ring with $n$ sections (which we call an $n$-ring), so the answer is given by $f(6)$. For $n=2$, we compute $f(n)=4\cdot3=12$. For $n \ge 3$, we can count the number of valid colorings as follows: choose one of the sections arbitrarily, which we may color in $4$ ways. Moving clockwise around the ring, there are $3$ ways to color each of the $n-1$ other sections. Therefore, we have $4 \cdot 3^{n-1}$ colorings of an $n$-ring.

 However, note that the first and last sections could be the same color under this method. To count these invalid colorings, we see that by "merging" the first and last sections into one, we get a valid coloring of an $(n-1)$-ring. That is, there are $f(n-1)$ colorings of an $n$-ring in which the first and last sections have the same color. Thus, $f(n) = 4 \cdot 3^{n-1} - f(n-1)$ for all $n \ge 3$.

 To compute the requested value $f(6)$, we repeatedly apply this formula: 
 \begin{align*} f(6)&=4\cdot3^5-f(5)\\&=4\cdot3^5-4\cdot3^4+f(4)\\&=4\cdot3^5-4\cdot3^4+4\cdot3^3-f(3)\\&=4\cdot3^5-4\cdot3^4+4\cdot3^3-4\cdot3^2+f(2)\\&=4(3^5-3^4+3^3-3^2+3)\\&=4\cdot3\cdot\frac{3^5+1}{3+1}\\&=\boxed{732}.\end{align*}
 (Solution by Vedantu .)

 





\textbf{Solution 4b (solution for all values of n)}

Do the same as solution 4 to get $f(n) = 4 \cdot 3^{n-1} - f(n-1)$. Then, multiply both sides by $(-1)^{n}$ to get $(-1)^{n}f(n) = (-1)^{n-1}f(n-1) - 4 \cdot (-3)^{n-1}$. Let $g(n) = (-1)^{n}f(n)$, $g(n) = g(n-1) - 4 \cdot (-3)^{n-1}$. $g(n) = g(2) - 4((-3)^{2}+(-3)^{3}+ \cdots + (-3)^{n-1}) = g(2) - 9(1-(-3)^{n-2})$. $f(n) = (-1)^{n}g(n) = (-1)^{n}g(2)-9(-1)^{n}+3^{n} = \boxed{3(-1)^{n}+3^{n}}$.

 





\textbf{Solution 5}

Label the sections 1, 2, 3, 4, 5, 6 clockwise. We do casework on the colors of sections 1, 3, 5.

 Case 1: the colors of the three sections are the same.In this case, each of sections 2, 4, 6 can be one of 3 colors, so this case yields $4 \times 3^3 = 108$ ways.

 Case 2: two of sections 1, 3, 5 are the same color.Note that there are 3 ways for which two of the three sections have the same color, and $4 \times 3 = 12$ ways to determine their colors. After this, the section between the two with the same color can be one of 3 colors and the other two can be one of 2 colors. This case yields $3 \times (4 \times 3) \times (3 \times 2 \times 2) = 432$ ways.

 Case 3: all three sections of 1, 3, 5 are of different colors.Clearly, there are $4 \times 3 \times 2 = 24$ choices for which three colors to use, and there are 2 ways to choose the colors of each of sections 2, 4, 6. Thus, this case gives $4 \times 3 \times 2 \times 2^3 = 192$ ways.

 In total, there are $108 + 432 + 192 = \boxed{732}$ valid colorings.

 Solution by ADMathNoob

 





\textbf{Solution 6}

We will take a recursive approach to this problem. We can start by writing out the number of colorings for a circle with $1, 2,$ and $3$ compartments, which are $4, 12,$ and $24.$ Now we will try to find a recursive formula, $C(n)$, for a circle with an arbitrary number of compartments $n.$ We will do this by focusing on the $n-1$ section in the circle. This section can either be the same color as the first compartment, or it can be a different color as the first compartment. We will focus on each case separately.

 Case 1:

 If they are the same color, we can say there are $C(n-2)$ to fill the first $n-1$ compartments. The $nth$ compartment must be different from the first and second to last compartments, which are the same color. Hence this case adds $3*C(n-2)$ to our recursive formula.

 Case 2:

 If they are different colors, we can say that there are $C(n-1)$ to fill the first $n-1$ compartments, and for the the $nth$ compartment, there are $2$ ways to color it because the $n-1$ and $1$ compartments are different colors. Hence this case adds $2*C(n-1).$

 So our recursive formula, $C(n)$, is $3*C(n-2) + 2*C(n-1).$ Using the initial values we calculated, we can evaluate this recursive formula up to $n=6.$ When $n=6,$ we get $\boxed{732}$ valid colorings.

 Solution by NeeNeeMath

 





\textbf{Solution 7}

WLOG, color the top left section $1$ and the top right section $2$. Then the left section can be colored $2$, $3$, or $4$, and the right section can be colored $1$, $3$, or $4$. There are $3 \cdot 3 = 9$ ways to color the left and right sections. We split this up into two cases.

 Case 1: The left and right sections are of the same color. There are $2$ ways this can happen: either they both are $3$ or they both are $4$. We have $3$ colors to choose for the bottom left, and $2$ remaining colors to choose for the bottom right, for a total of $2 \cdot 3 \cdot 2$ cases.

 Case 2: The left and right sections are of different colors. There are $9 - 2 = 7$ ways this can happen. Assume the left is $3$ and the right is $4$. Then the bottom left can be $1$, $2$, or $4$, and the bottom right can be $1$, $2$, or $3$. However the bottom sections cannot both be $1$ or both be $2$, so there are $3 \cdot 3 - 2 = 7$ ways to color the bottom sections, for a total for $7 \cdot 7 = 49$ colorings.

 Since there were $4 \cdot 3 = 12$ ways to color the top sections, the answer is $12 \cdot (12 + 49) = \boxed{732}$.

 





\textbf{Solution 8}

Label the four colors $1, 2, 3,$ and $4$, respectively.Now let's imagine a circle with the four numbers $1, 2, 3,$ and $4$ written clockwise. We'll say that a bug is standing on number $a$. It is easy to see that for the bug to move to a different number, it must walk $1, 2,$ or $3$ steps clockwise. (This is since adjacent numbers can't be the same, as stated in the problem). Note that the sixth number in the bug's walking sequence must not equal the first number. Thus, our total number of ways, $N$, given the bug's starting number $k$, is simply the number of ordered quintuplets of positive integers $(a_1, a_2, a_3, a_4, a_5)$ that satisfy $a_i \in \{1, 2, 3\}$ for all $1 \leq i \leq 5$ and
 \[\sum_{i=1}^{5} a_i \neq 8, 12,\]since the bug cannot land on $k$ again on his fifth and last step.We know that the number of ordered quintuplets of positive integers $(a_1, a_2, a_3, a_4, a_5)$ that satisfy $a_i \in \{1, 2, 3\}$ without the other restriction is just $3^5$, so we aim to find the number of quintuplets such that
 \[a_1 + a_2 + a_3 + a_4 + a_5 = 8, 12.\]Note that the number of ordered quintuplets satisfying $\sum_{i=1}^{5} a_i = 8$ is the same as the number of them satisfying $\sum_{i=1}^{5} a_i = 12$ due to symmetry. By stars and bars, there are $\dbinom{7}{4} = 35$ ways to distribute the three "extra" units to the five variables $a_1, a_2, a_3, a_4, a_5$ (since $a_i \geq 1$), but ways of distribution such that one variable is equal to $4$ are illegal, so the actual number of ways is $35 - 5 = 30$. Since there are four possible values of $k$ (or the starting position for the bug), we obtain
 \[N = 4 \times (3^5 - 2 \times 30) = \boxed{732}.\]-fidgetboss\_4000

 





\textbf{Solution 9}

Let's number the regions $1,2,\dots 6$. Suppose we color regions $1,2,3$. Then, how many ways are there to color $4,5,6$?

 Note: the numbers are numbered as shown:

 
                    \begin{figure}[H]
                        \centering
                        \includegraphics[width=0.5\textwidth]{Solution/AIME_2016_II_12_Solution_Figure_1.png}
                    \end{figure}
                    

 $\textbf{Case 1:}$ The colors of $1,2,3$ are $BAB$, in that order.

 Then the colors of $6,5,4$ can be $AXA$, $AXC$, $CXA$, $CXC$, or $CXD$ in that order, where $X$ is any color not equal to its surroundings. Then there are $4$ choices for $A$, $3$ choices for $B$ (it cannot be $A$), $2$ choices for $C$, and $1$ for $D$, the last color. So, summing up, we have 
 \[4*3(1*3+2*2+2*2+2*3+2*1*2)=12*21=252\] colorings.

 $\textbf{Case 2:}$ The colors of $1,2,3$ are $BAC$, in that order.

 Again, we list out the possible arrangements of $6,5,4$: $AXA$, $AXB$, $CXA$, $AXD$, $DXA$, $CXB$, $CXD$, $DXB$, or $DXD$. (Easily simplified; listed here for clarity.) Then there are $4$ choices for $A$ as usual, $3$ choices for $B$, and so on. Hence we have 
 \[4*3*2(1*3+1*2+1*2+1*2+1*2+1*2+1*2+1*2+1*3)=24*20=480\] colorings in this case.

 Adding up, we have $252+480=\boxed{732}$ as our answer.

 This solution was brought to you by LEONARD\_MY\_DUDE

 





\textbf{Solution 10 (chromatic polynomial-VERY HELPFUL)}

We quickly notice that this is just the cycle graph with 6 vertices. The chromatic polynomial for a cycle is $(k-1)^n+(-1)^n (k-1)$ where we use $k$ colors on a cycle of $n$ vertices. Plugging in $k=4$ and $n=6$ we arrive at $\boxed{732}$.~chrisdiamond10

 





\textbf{Solution 11 (step-by-step case analysis)}

Let's label the regions as $1,2,3,4,5,6$ in that order. We start with region $1$. There are no restrictions on the color of region $1$ so it can be any of the four colors. We know move on the region $2$. It can be any color but color used for region $1$, giving us $3$ choices. Section $3$ is where it gets a bit complicated; we will have to do casework based on whether the color of region $3$ is that of region $1$ or not.

 If we have the color of region $3$ being different from that of region $1$ (in which we color do so in $2$ ways), then we have need for another casework at region $4$. If the color of region $4$ is different from that of region $1$ (which can be achieved in $2$ ways), then we have yet another casework split.

 If the color of region $5$ is different from that of region $1$ (which can be done in $2$ ways, then we would have a total of $2$ possible colorings for region $6$ (for it cannot be the color of regions $1$ nor $5$). Moving on to the case where the color of section $5$ is the same as that of section $1$ (which can be done in $1$ way), we will have $3$ ways (region $6$ cannot be that color of both region $1$ and $5$). Thus, if the color of region $4$ is different of that of region $1$, then we have $2\cdot 2 + 3 = 7$ ways.

 If the color of region $4$ is the same as that of region $1$ (which can be done in one way), then the color of region $5$ have to be different from that of sector $1$ ($3$ ways). That means there will be $2$ choices for the color of region $6$. So if the color of region $4$ is the same as region $1$, then we will have $3\cdot 2 = 6$. That means if the color of section $3$ is different from that of region $1$, then there are $2\cdot 7 + 6 = 20$.

 Now, moving on to the case where section $3$ has the same color of region $1$. That means section $4$ will have to be a different color of that of region $1$ ($3$ ways). So, that means we have region $5$ to be either the same color or different color as region $1$. If it is different (which can be done in $2$ ways), then there will be $2$ possibilities on the color of sector $6$. If it is same (which can be done in $1$ way), then there are $3$ ways to color region $6$. So, if section $3$ has the same color as section $1$, then we have $3(2\cdot 2 + 3) = 3(7) = 21$.

 Now, in overall we will have $4\cdot 3 (2(20)+21) = 12(61) = \boxed{732}$.

 The full expansion without the explanation is right here:

 
 \[4\cdot 3(2(2\cdot 7 + 3\cdot 2) + 3(2\cdot 2 + 3)) = 732\]

 





\textbf{Solution 12 (Probability & Expected Value)}

Let the top-right segment be segment $1,$ and remaining segments are numbered $2,3,4,5,6$ in clockwise order.We have $4$ choices for segment $1,$ and $3$ choices for segments $2,3,4,5.$ For segment $6,$ we wish to find the expected value of the number of choices for segment $6$'s color, which depends on whether segment $5$ is red. If we let $P(n)$ denote the probability of segment $P$ being the same color as segment $1$ (for simplicity, denote segment $1$'s color as red), we get the following recurrence relation:
 \[P(n)=\dfrac{1-P(n-1)}{3}.\]This is because that you cannot have two reds in a row (hence the $1-P(n-1)$) and if segment $n-1$ is not red, there are three possible colors, one of which is red (hence the divide by $3$). Using the obvious fact that $P(1)=1$ by the definition of $P(n),$ we find that
 \[P(5)=\dfrac{7}{27}.\]If segment $5$ is red, then there are three possible colors for segment $6.$ If it is not, there are $2$ possible choices for segment  $6.$ This means the expected number of color choices for segment $6$ is
 \[\dfrac{7}{27}\cdot3+\dfrac{20}{27}\cdot2,\]and the total number of colorings of the ring is 
 \[4\cdot3\cdot3\cdot3\cdot3\cdot\left(\dfrac{7}{27}\cdot3+\dfrac{20}{27}\cdot2\right)=\boxed{732}.\]

 




	\item (\textit{AIME II 2016}) Beatrix is going to place six rooks on a $6 \times 6$ chessboard where both the rows and columns are labeled $1$ to $6$; the rooks are placed so that no two rooks are in the same row or the same column. The \textit{value} of a square is the sum of its row number and column number. The \textit{score} of an arrangement of rooks is the least value of any occupied square. The average score over all valid configurations is $\frac{p}{q}$, where $p$ and $q$ are relatively prime positive integers. Find $p+q$.


 
 \textbf{Answer: }371



\textbf{Solution 1}

We casework to find the number of ways to get each possible score. Note that the lowest possible score is $2$ and the highest possible score is $7$. Let the bijective function $f(x)=\{1,2,3,4,5,6\} \to \{1,2,3,4,5,6\}$ denote the row number of the rook for the corresponding column number.

 
                \begin{itemize}
                    \item For a score of $2$, we must have $f(1)=1$, and we can arrange the rest of the function however we want, so there are $5!=120$ ways.
                \end{itemize}


                
                \begin{itemize}
                    \item For a score of $3$, we must have either $f(1)=2$ or $f(2)=1$, and we can arrange the rest of the rooks however we want, so by PIE the number of ways is $5!+5!-4!=216$.
                \end{itemize}


                
                \begin{itemize}
                    \item For a score of $4$, we must have $f(1)=3$, $f(2)=2$, or $f(3)=1$. If $f(1)=3$, we just don't want $f(2)=1$, if $f(2)=2$, we don't want $f(1)=1$, or if $f(3)=1$, we don't want $f(1)=2$, otherwise we can arrange the function however we like. If at least $2$ of the values rooks have a value of $4$, we can arrange the rest of the rooks however we like, so there are $3(5!-4!)-3(4!)+3!=222$ by PIE.
                \end{itemize}


                
                \begin{itemize}
                    \item If the score is $5$, then we have either $f(4)=1$, $f(3)=2$, $f(2)=3$, or $f(1)=4$. If we have the first case, we don't want $f(2)=2$, $f(1)=2$, or $f(1)=3$, so by PIE the number of bad cases is $4!+4!+4!-3!=66$. If we have the second case, then we don't want $f(2)=1$, $f(1)=1$, or $f(1)=3$, so similarly there are $66$ bad cases. This is true for the other two cases as well. Therefore, there are a total of $120-66=54$ good cases for each one. The number of ways to get $f(4)=1, f(1)=4$ is $4!-3!=18$ because we don't want $f(2)=2$, the number of ways to get $f(4)=1, f(2)=3$ is $18$ ways because we don't want $f(1)=2$, the number of ways to get $f(4)=1, f(3)=2$ is $18$ ways because we don't want $f(1)=3$, the number of ways to get $f(3)=2, f(1)=4$ is $18$ ways because we don't want $f(2)=1$, the number of ways to get $f(3)=2, f(2)=3$ is $18$ ways because we don't want $f(1)=1$, and the number of ways to get $f(1)=4, f(2)=3$ is $18$ ways because we don't want $f(3)=1$. The number of ways to get at least $3$ cases satisfied is $3! \cdot 4$ because we can arrange the remaining rooks however we like, and the number of ways to get all $4$ cases satisfied is $2!=2$  ways because we can arrange the remaining rooks however we like, so by PIE we have $54 \cdot 4-18 \cdot 6 + 3! \cdot 4-2!=130$ ways to get a score of $5$.
                \end{itemize}


                
                \begin{itemize}
                    \item The only way to get a score of $7$ is to have all the rooks run on the antidiagonal. Therefore, the number of ways to get a sum of $6$ is $6!-120-216-222-130-1=31$.
                \end{itemize}


                Thus, the expected sum is $\dfrac{120 \cdot 2 + 216 \cdot 3 + 222 \cdot 4 + 130 \cdot 5 + 31 \cdot 6 + 1 \cdot 7}{720}= \dfrac{2619}{720}=\dfrac{291}{80}$, so the desired answer is $291+80=\boxed{371}$.

 





\textbf{Solution 2}

If the score is $n+1$, then one of the rooks must appear in the $n$th antidiagonal, and this is the first antidiagonal in which a rook can appear. To demonstrate this, we draw the following diagram when $n=4$.

 
                    \begin{figure}[H]
                        \centering
                        \includegraphics[width=0.5\textwidth]{Solution/AIME_2016_II_13_Solution_Figure_1.png}
                    \end{figure}
                    We first count the number of arrangements that avoid the squares above the $n$th diagonal, and then we subtract from these the number of arrangements that avoid all squares above the $(n+1)$th diagonal. In the first column, there are $7-n$ rows in which to place the rook. In the second column, there is one more possible row, but one of the rows is used up by the rook in the first column, hence there are still $7-n$ places to place the rook. This pattern continues through the $n$th column, so there are $(7-n)^n$ ways to place the first $n$ rooks while avoiding the crossed out squares. We can similarly compute that there are $(6-n)^n$ ways to place the rooks in the first $n$ columns that avoid both the crossed out and shaded squares. Therefore, there are $(7-n)^n-(6-n)^n$ ways to place the first $n$ rooks such that at least one of them appears in a shaded square.

 After this, there are $(6-n)$ rows and $(6-n)$ columns in which to place the remaining rooks, and we can do this in $(6-n)!$ ways. Hence the number of arrangements with a score of $n$ is $((7-n)^n-(6-n)^n)\cdot(6-n)!$. We also know that $n$ can range from from $1$ to $6$, so the average score is given by

 
 \[\frac{2\cdot(6^1-5^1)\cdot5!+3\cdot(5^2-4^2)\cdot4!+4\cdot(4^3-3^3)\cdot3!+5\cdot(3^4-2^4)\cdot 2!+6\cdot(2^5-1^5)\cdot 1!+7\cdot(1^6-0^6)\cdot 0!}{6!}=\frac{291}{80}.\]Thus the answer is $\boxed{371}$.

 


 





\textbf{Solution 3}

So we first count the number of permutations with score $\ge 2$. This is obviously $6!=720$. Then, the number of permutations with score $\ge 3$ can also be computed: in the first column, there are five ways to place a rook- anywhere but the place with score $1$. In the next column, there are $5$ ways to place a rook- anywhere but the one in the same row as the previous row. We can continue this to obtain that the number of permutations with score $\ge 3$ is $600$. Doing the same for scores $\ge 4$, $\ge 5$, $\ge 6$, and $\ge 7$ we obtain that these respective numbers are $384$, $162$, $32$, $1$. 

 Now, note that if $a_k$ is the number of permutations with score $\ge k$, then $a_k-a_{k-1}=b_{k}$, where $b_k$ is the number of permutations with score exactly $k$. Thus, we can compute the number of permutations with scores $2$, $3$, etc as $120,216,222,130,31,1$. We then compute 
 \[\frac{120(2)+216(3)+222(4)+130(5)+31(6)+1(7)}{720}=\frac{291}{80}\] leading us to the answer of $291+80=\boxed{371}$. $\blacksquare$

 \textbf{Note} This problem is pretty bashy. There isn't a clever way to solve it. 

 


 





\textbf{Solution 4 (builds off Solution 3)}

The problem is asking us to compute $\mathbb{E}[S]$, where $S$ is the random variable that takes an arrangement of rooks and outputs its score, which is a non-negative integer quantity. For any random variable $S$ with non-negative integer values, we have the tail sum formula
 \[\mathbb{E}[S] = \sum_{n = 1}^{\infty}\mathbb{P}(S\geq n).\]These probabilities can be computed as in Solution 3, giving us the following table.

 




	\item (\textit{AIME II 2016}) Equilateral $\triangle ABC$ has side length $600$. Points $P$ and $Q$ lie outside the plane of $\triangle ABC$ and are on opposite sides of the plane. Furthermore, $PA=PB=PC$, and $QA=QB=QC$, and the planes of $\triangle PAB$ and $\triangle QAB$ form a $120^{\circ}$ dihedral angle (the angle between the two planes). There is a point $O$ whose distance from each of $A,B,C,P,$ and $Q$ is $d$. Find $d$.


 
 \textbf{Answer: }450



\textbf{Solution 1}

The inradius of $\triangle ABC$ is $100\sqrt 3$ and the circumradius is $200 \sqrt 3$. Now, consider the line perpendicular to plane $ABC$ through the circumcenter of $\triangle ABC$. Note that $P,Q,O$ must lie on that line to be equidistant from each of the triangle's vertices. Also, note that since $P, Q, O$ are collinear, and $OP=OQ$, we must have $O$ is the midpoint of $PQ$.  Now, Let $K$ be the circumcenter of $\triangle ABC$, and $L$ be the foot of the altitude from $C$ to $AB$. We must have $\tan(\angle KLP+ \angle QLK)= \tan(120^{\circ})$. Setting $KP=x$ and $KQ=y$, assuming WLOG $x>y$, we must have $\tan(120^{\circ})=-\sqrt{3}=\dfrac{\dfrac{x+y}{100 \sqrt{3}}}{\dfrac{30000-xy}{30000}}$. Therefore, we must have $100(x+y)=xy-30000$. Also, we must have $\left(\dfrac{x+y}{2}\right)^{2}=\left(\dfrac{x-y}{2}\right)^{2}+120000$ by the Pythagorean theorem, so we have $xy=120000$, so substituting into the other equation we have $90000=100(x+y)$, or $x+y=900$. Since we want $\dfrac{x+y}{2}$, the desired answer is $\boxed{450}$.

 





\textbf{Solution 2 (Short & Simple)}

Draw a good diagram. Draw $CH$ as an altitude of the triangle. Scale everything down by a factor of $100\sqrt{3}$, so that $AB=2\sqrt{3}$. Finally, call the center of the triangle U. Draw a cross-section of the triangle via line $CH$, which of course includes $P, Q$. From there, we can call $OU=h$. There are two crucial equations we can thus generate. WLOG set $PU<QU$, then we call $PU=d-h, QU=d+h$. First equation: using the Pythagorean Theorem on $\triangle UOB$, $h^2+2^2=d^2$. Next, using the tangent addition formula on angles $\angle PHU, \angle UHQ$ we see that after simplifying $-d^2+h^2=-4, 2d=3\sqrt{3}$ in the numerator, so $d=\frac{3\sqrt{3}}{2}$. Multiply back the scalar and you get $\boxed{450}$. Not that hard, was it?

 





\textbf{Solution 3}

To make numbers more feasible, we'll scale everything down by a factor of $100$ so that $\overline{AB}=\overline{BC}=\overline{AC}=6$. We should also note that $P$ and $Q$ must lie on the line that is perpendicular to the plane of $ABC$ and also passes through the circumcenter of $ABC$ (due to $P$ and $Q$ being equidistant from $A$, $B$, $C$), let $D$ be the altitude from $C$ to $AB$. We can draw a vertical cross-section of the figure then: 
                    \begin{figure}[H]
                        \centering
                        \includegraphics[width=0.5\textwidth]{Solution/AIME_2016_II_14_Solution_Figure_1.png}
                    \end{figure}
                     We let $\angle PDI=\alpha$ so $\angle QDI=120^{\circ}-\alpha$, also note that $\overline{PO}=\overline{QO}=\overline{CO}=d$. Because $I$ is the centroid of $ABC$, we know that ratio of $\overline{CI}$ to $\overline{DI}$ is $2:1$. Since we've scaled the figure down, the length of $CD$ is $3\sqrt{3}$, from this it's easy to know that $\overline{CI}=2\sqrt{3}$ and $\overline{DI}=\sqrt{3}$. The following two equations arise: 
 \begin{align} \sqrt{3}\tan{\left(\alpha\right)}+\sqrt{3}\tan{\left(120^{\circ}-\alpha\right)}&=2d \\ \sqrt{3}\tan{\left(\alpha\right)} - d &= \sqrt{d^{2}-12} \end{align} Using trig identities for the tangent, we find that 
 \begin{align*} \sqrt{3}\tan{\left(120^{\circ}-\alpha\right)}&=\sqrt{3}\left(\frac{\tan{\left(120^{\circ}\right)}+\tan{\left(\text{-}\alpha\right)}}{1-\tan{\left(120^{\circ}\right)}\tan{\left(\text{-}\alpha\right)}}\right) \\ &= \sqrt{3}\left(\frac{\text{-}\sqrt{3}+\tan{\left(\text{-}\alpha\right)}}{1+\sqrt{3}\tan{\left(\text{-}\alpha\right)}}\right) \\ &= \sqrt{3}\left(\frac{\text{-}\sqrt{3}-\tan{\left(\alpha\right)}}{1-\sqrt{3}\tan{\left(\alpha\right)}}\right) \\ &= \frac{\sqrt{3}\tan{\left(\alpha\right)}+3}{\sqrt{3}\tan{\left(\alpha\right)}-1}.\end{align*}
 Okay, now we can plug this into $\left(1\right)$ to get: 
 \begin{align}\sqrt{3}\tan{\left(\alpha\right)}+\frac{\sqrt{3}\tan{\left(\alpha\right)}+3}{\sqrt{3}\tan{\left(\alpha\right)}-1}&=2d \\ \sqrt{3}\tan{\left(\alpha\right)} - d &= \sqrt{d^{2}-12} \end{align} Notice that $\alpha$ only appears in the above system of equations in the form of $\sqrt{3}\tan{\left(\alpha\right)}$, we can set $\sqrt{3}\tan{\left(\alpha\right)}=a$ for convenience since we really only care about $d$. Now we have 
 \begin{align}a+\frac{a+3}{a-1}&=2d \\ a - d &= \sqrt{d^{2}-12} \end{align} Looking at $\left(2\right)$, it's tempting to square it to get rid of the square-root so now we have: 
 \begin{align*}a^{2}-2ad+d^{2}&=d^{2}-12 \\ a - 2ad &= \text{-}12 \end{align*}
 See the sneaky $2d$ in the above equation? That we means we can substitute it for $a+\frac{a+3}{a-1}$: 
 \begin{align*}a^{2}-2ad+d^{2}&=d^{2}-12 \\ a^{2} - a\left(a+\frac{a+3}{a-1}\right) &= \text{-}12 \\ a^{2}-a^{2}-\frac{a^{2}+3a}{a-1} &=\text{-}12 \\ -\frac{a^{2}+3a}{a-1}&=\text{-}12 \\ \text{-}a^{2}-3a&=\text{-}12a+12 \\ 0 &= a^{2}-9a+12 \end{align*}
 Use the quadratic formula, we find that $a=\frac{9\pm\sqrt{9^{2}-4\left(1\right)\left(12\right)}}{2\left(1\right)}=\frac{9\pm\sqrt{33}}{2}$ - the two solutions were expected because $a$ can be $\angle PDI$ or $\angle QDI$. We can plug this into $\left(1\right)$: 
 \begin{align*}a+\frac{a+3}{a-1}&=2d \\ \frac{9\pm\sqrt{33}}{2}+\frac{\frac{9\pm\sqrt{33}}{2}+3}{\frac{9\pm\sqrt{33}}{2}-1}=2d \\ \frac{9\pm\sqrt{33}}{2}+\frac{15\pm\sqrt{33}}{7\pm\sqrt{33}} &= 2d\end{align*}
 I'll use $a=\frac{9+\sqrt{33}}{2}$ because both values should give the same answer for $d$. 
 \begin{align*} \frac{9+\sqrt{33}}{2}+\frac{15+\sqrt{33}}{7+\sqrt{33}} &= 2d \\ \frac{\left(9+\sqrt{33}\right)\left(7+\sqrt{33}\right)+\left(2\right)\left(15+\sqrt{33}\right)}{\left(2\right)\left(7+\sqrt{33}\right)} &= 2d \\ \frac{63+33+16\sqrt{33}+30+2\sqrt{33}}{14+2\sqrt{33}} &= 2d \\ \frac{126+18\sqrt{33}}{14+2\sqrt{33}} &= 2d \\ 9 &= 2d \\ \frac{9}{2} &= d\end{align*}
 Wait! Before you get excited, remember that we scaled the entire figure by $100$?? That means that the answer is $d=100\times\frac{9}{2}=\boxed{450}$.-fatan

 bro i would be horrified if i saw a fraction as my answer after this much work on the aime wdym “excited”

 





\textbf{Solution 4}

We use the diagram from solution 3. From basic angle chasing,
 \[180=\angle{QOC}+\angle{CO}P=2\angle{OCP}+2\angle{OCQ}=2\angle{QCP}\]so triangle QCP is a right triangle. This means that triangles $CQI$ and $CPI$ are similar. If we let $\angle{IDQ}=x$ and $\angle{PDI}=y$, then we know $x+y=120$ and 
 \[\frac{PI}{IC}=\frac{IC}{IQ}\Rightarrow\frac{100\sqrt{3}\tan{y}}{200\sqrt{3}}=\frac{200\sqrt{3}}{100\sqrt{3}\tan{x}}\Rightarrow\tan{x}\tan{y}=4\] We also know that 
 \[PQ=2d=100\sqrt{3}(\tan{x}+\tan{y})\] 
 \[d=50\sqrt{3}(\tan{x}+\tan{y})\] 
 \[\frac{d}{1-\tan{x}\tan{y}}=50\sqrt{3}\cdot\frac{\tan{x}+\tan{y}}{1-\tan{x}\tan{y}}\] 
 \[\frac{d}{-3}=50\sqrt{3}\tan{(x+y)}\] 
 \[d=-150\sqrt{3}\tan{120}=-150\sqrt{3}(-\sqrt{3})=\boxed{450}\]

 -EZmath2006

 





\textbf{Solution 5}

We use the diagram from solution 3.

 Let $BP = a$ and $BQ = b$. Then, by Stewart's on $BPQ$, we find
 \[2x^3 + 2x^3 = a^2x + b^2x \implies a^2 + b^2 = 4x^2.\]

 The altitude from $P$ to $ABC$ is $\sqrt{a^2 - (200\sqrt{3})^2}$ so
 \[PQ = 2x = \sqrt{a^2 - (200\sqrt{3})^2} + \sqrt{b^2 - (200\sqrt{3})^2}.\]

 Furthermore, the altitude from $P$ to $AB$ is $\sqrt{a^2 - 300^2}$, so, by LoC and the dihedral condition,
 \[a^2 - 300^2 + b^2 - 300^2 + \sqrt{a^2 - 300^2}\sqrt{b^2-300^2} = 4x^2.\]

 Squaring the equation for $PQ$ and substituting $a^2 + b^2  = 4x^2$ yields
 \[2\sqrt{a^2 - (200\sqrt{3})^2}\sqrt{b^2 - (200\sqrt{3})^2} = 6\cdot 200^2.\]

 Substituting $a^2 + b^2 = 4x^2$ into the other equation,
 \[\sqrt{a^2 - 300^2}\sqrt{b^2-300^2} = 2\cdot 300^2.\]

 Squaring both of these gives
 \[a^2b^2-3\cdot 200^2(a^2 + b^2) + 9\cdot 200^4 = 9\cdot 200^4\]
 \[a^2b^2 - 300^2(a^2+b^2) + 300^4 = 4\cdot 300^4.\]

 Substituting $a^2 + b^2 = 4x^2$ and solving for $x$ gives $\boxed{450}$, as desired.

 -mathtiger6

 





\textbf{Solution 6 (Geometry)}


                    \begin{figure}[H]
                        \centering
                        \includegraphics[width=0.5\textwidth]{Solution/AIME_2016_II_14_Solution_Figure_2.png}
                    \end{figure}
                    
                    \begin{figure}[H]
                        \centering
                        \includegraphics[width=0.5\textwidth]{Solution/AIME_2016_II_14_Solution_Figure_3.png}
                    \end{figure}
                    Let $AB = a, M$ be midpoint $BC, I$ be the center of equilateral $\triangle ABC,$$IM = b = \frac {a}{2\sqrt{3}}, O$ be the center of sphere $ABCPQ.$Then 
 \[AI = 2b, AO = BO = PO =QO = d.\]
 \[QA=QB=QC,PA=PB=PC \implies\]
 \[POIQ\perp ABC, \angle PMQ = 120^\circ.\](See upper diagram).

 We construct the circle PQMD, use the formulas for  intersecting chords and get
 \[DI = 5b, FI = EO = \frac{3b}{2}\]
 \[\implies FM = \frac{5b}{2}.\](See lower diagram).

 We apply the Law of Sine to $\triangle PMQ$ and get
 \[2EM \sin 120^\circ =PQ\]
 \[\implies r \sqrt{3} = 2d\]
 \[\implies 3r^2 = 4d^2.\] We apply the Pythagorean Law on $\triangle AOI$ and $\triangle EFM$ and get
 \[d^2 = 4b^2 + OI^2, r^2 = \frac {25b^2}{4} + EF^2 \implies\]
 \[r = 3b\implies d = \frac {3a}{2} = \boxed {450}.\]\textbf{vladimir.shelomovskii@gmail.com, vvsss}

 





\textbf{Solution 7}

Let $M$ be the midpoint of $\overline{AB}$ and $X$ the center of $\triangle ABC$. Then 
 \[P, O, Q, M, X, C\] all lie in the same vertical plane. We can make the following observations:

 
                \begin{itemize}
                    \item The equilateral triangle has side length $600$, so $MC=300\sqrt{3}$ and $X$ divides $MC$ so that $MX=100\sqrt{3}$ and $XC=200\sqrt{3}$;\item $O$ is the midpoint of $PQ$ since $O$ is equidistant from $A, B, C, P, Q$ – it is also the circumcenter of $\triangle PCQ$;\item $\angle PMQ=120^{\circ}$, the dihedral angle.
                \end{itemize}


                To make calculations easier, we will denote $100\sqrt{3}=m$, so that $MX=m$ and $XC=2m$.

 
                    \begin{figure}[H]
                        \centering
                        \includegraphics[width=0.5\textwidth]{Solution/AIME_2016_II_14_Solution_Figure_4.png}
                    \end{figure}
                    

 Denote $PX=p$ and $QX=q$, where the tangent addition formula on $\triangle PMQ$ yields 
 \[\frac{\tan\measuredangle PMX+\tan\measuredangle QMX}{1-\tan\measuredangle PMX\tan\measuredangle QMX}=\tan(120^{\circ})=-\sqrt{3}.\] Using $\tan\measuredangle PMX=\frac{p}{m}$ and $\tan\measuredangle QMX=\frac{q}{m}$, we have 
 \[\frac{\frac{p}{m}+\frac{q}{m}}{1-\frac{p}{m}\cdot\frac{q}{m}}=-\sqrt{3}.\] After multiplying both numerator and denominator by $m^{2}$ we have 
 \[\frac{(p+q)m}{m^{2}-pq}=-\sqrt{3}.\] But note that $pq=(2m)(2m)=4m^{2}$ by power of a point at $X$, where we deduce by symmetry that $MM^{\prime}=MX=m$ on the diagram below: 
                    \begin{figure}[H]
                        \centering
                        \includegraphics[width=0.5\textwidth]{Solution/AIME_2016_II_14_Solution_Figure_5.png}
                    \end{figure}
                    

 Thus 
 \begin{align*} \frac{(p+q)m}{m^{2}-4m^{2}}=-\sqrt{3} \\ \frac{p+q}{-3m}=-\sqrt{3} \\ p+q=\left(-\sqrt{3}\right)\left(-3m\right) \\ p+q=3\sqrt{3}\cdot m.\end{align*}
 Earlier we assigned the variable $m$ to the length $100\sqrt{3}$ which implies $PQ=\left(3\sqrt{3}\right)\left(100\sqrt{3}\right)=900$. Thus the distance $d$ is equal to $\frac{PQ}{2}=\boxed{450}$.

 





\textbf{Solution 8 (Law of Cosines)}

Let $Z$ be the center of $\triangle ABC$. Let $A’$ be the midpoint of $AB$. Let $ZA’ = c = 100\sqrt{3}$ and $ZA = 2c = 200\sqrt{3}$. Let $PZ = a$ and $QZ = b$. We will be working in the plane that contains the points: $A$, $P$, $A’$, $Q$, $O$, and $Z$.

 Since $P$, $O$, and $Q$ are collinear and $PO = QO = AO$, $\triangle PAQ$ is a right triangle with $\angle PAQ = 90^{\circ}$. Since $AZ \perp PQ$, $(PZ)(QZ) = (AZ)^2 = ab = (2c)^2 = 120000$.

 $PA’ = \sqrt{a^2 + c^2}$, $QA’ = \sqrt{b^2 + c^2}$, $PQ = a + b$, and $\angle PA'Q = 120^{\circ}$. By Law of Cosines 
 \[(a + b)^2 = a^2 + b^2 + 2c^2 + \sqrt{a^2b^2 + a^2c^2 + b^2c^2 + c^4}\]. Substituting $4c^2$ for $ab$ and simplifying, we get 
 \[6c = \sqrt{17c^2 + a^2 + b^2}\]. Squaring and simplifying, we get 
 \[a^2 + b^2 = 19c^2 = 570000\]. Adding $2ab = 8c^2$ to both sides we get $PQ = a + b = 900$. Since $O$ is the midpoint of $PQ$, $d = PO = \boxed{450}$

 





\textbf{Solution 9 (Coordinate Bash)}

Set $AB = s, M$ as midpoint $BC, I$ as the center of equilateral $\triangle ABC,$ and by $30-60-90$ triangle formulas, we know $IM = b = \frac {s}{2\sqrt{3}}, IA = IB = IC = \frac{s}{\sqrt{3}}$.

 First note that $P$ and $Q$ are on the line perpendicular to plane $ABC$ through the circumcenter of $\triangle ABC$. Then notice that point $O$ is the midpoint of $PQ$, so it also lies on this axis. Then, we see that the position of $O$ relative to the triangle fully determines the positions of both $P$ and $Q$, so setting the coordinates of $O$ as $(0,0,h)$ where the origin is defined the circumcenter and $\triangle ABC$ lies in the $xy$-plane, we get: 
 \begin{align*} OA^2 = OB^2 = OC^2 = h^2+\frac{s^2}{3} = OP^2 = OQ^2 \end{align*}


 Thus, the coordinates of $P$ and $Q$ are $(0,0,h\pm\sqrt{h^2+\frac{s^2}{3}})$. Now we make use of the angle condition. WLOG set the coordinates of $M$ as $(\frac{s}{2\sqrt{3}},0,0)$. We know the angle between $PM$ and $QM$ is $120^{\circ}$, so after solving for the vectors, taking their dot product, equating it to $PM \cdot QM \cdot\cos{120^{\circ}}$, and $\textit{finally}$ solving for $h$ in terms of $s$, we get 
 \[ h^2 = \frac{11}{48}s^2 \\ \implies OA^2 = h^2 + \frac{s^2}{3} = \frac{27}{48}s^2 \\ \implies OA = \frac{3}{4}s = \boxed{450} \]

 
(If someone feels like it, please feel free to fill in the rest of the details!)

 





\textbf{Solution 10 (Coordinate Geometry)}

We begin by scaling the entire diagram by a factor of $\frac{1}{50}$. We establish a coordinate system with $AC$ being the $x$-axis, where $A$ is the origin. Note that $P,Q,O$ all lie on the line $x=6, y=2\sqrt{3}$ in order to achieve the equidistant condition. Let $M$ be the midpoint of $AB$. Here is our setup:\begin{align*}A&=(0,0,0)\\B&=(6,6\sqrt{3},0)\\C&=(12,0,0)\\M&=(3,3\sqrt{3},0)\\P&=(6,2\sqrt{3},z_{1})\\Q&=(6,2\sqrt{3},z_{2})\\O&=(6,2\sqrt{3},\frac{z_{1}+z_{2}}{2})\end{align*}
Since $\overline{MP}\perp\overline{AB}$ and $\overline{MQ}\perp\overline{AB}$, $\angle PMQ$ is the dihedral angle formed by the planes of $\triangle PAB$ and $\triangle QAB$. By the distance formula and the Law of Cosine on $\triangle PMQ$, we obtain the following equations:\begin{align*}48+\frac{(z_1+z_2)^2}{4}&=\frac{(z_1-z_2)^2}{4}\\24+z_{1}^{2}+z_{2}^{2}+\sqrt{(z_{1}^{2}+12)(z_{2}^{2}+12)}&=(z_1-z_2)^2\end{align*}
We see that $z_1z_2=-48$, and the second equation simplifies to
 \[z_1^2z_2^2+32z_1z_2+144=4(z_1^2+z_2^2)\]This follows that $z_1^2+z_2^2=228$. Therefore,
 \[OP=\sqrt{\frac{(z_1-z_2)^2}{4}}=\sqrt{\frac{228+96}{4}}=9\]Scaling the distance by $50$ yields $d=9\cdot 50=\boxed{450}$, as desired.

 




	\item (\textit{AIME II 2016}) For $1 \leq i \leq 215$ let $a_i = \dfrac{1}{2^{i}}$ and $a_{216} = \dfrac{1}{2^{215}}$. Let $x_1, x_2, ..., x_{216}$ be positive real numbers such that $\sum_{i=1}^{216} x_i=1$ and $\sum_{1 \leq i < j \leq 216} x_ix_j = \dfrac{107}{215} + \sum_{i=1}^{216} \dfrac{a_i x_i^{2}}{2(1-a_i)}$. The maximum possible value of $x_2=\dfrac{m}{n}$, where $m$ and $n$ are relatively prime positive integers. Find $m+n$.


 
 \textbf{Answer: }863



\textbf{Solution 1}

Note that
 \begin{align*}\sum_{1 \leq i < j \leq 216} x_ix_j &= \frac12\left(\left(\sum_{i=1}^{216} x_i\right)^2-\sum_{i=1}^{216} x_i^2\right)\\&=\frac12\left(1-\sum x_i^2\right).\end{align*}
Substituting this into the second equation and collecting $x_i^2$ terms, we find 
 \[\sum_{i=1}^{216}\frac{x_i^2}{1-a_i}=\frac{1}{215}.\] Conveniently, $\sum_{i=1}^{216} 1-a_i=215$ so we find
 \[\left(\sum_{i=1}^{216} 1-a_i\right)\left(\sum_{i=1}^{216}\frac{x_i^2}{1-a_i}\right)=1=\left(\sum_{i=1}^{216} x_i\right)^2.\] This is the equality case of the Cauchy-Schwarz Inequality, so $x_i=c(1-a_i)$ for some constant $c$. Summing these equations and using the facts that $\sum_{i=1}^{216} x_i=1$ and $\sum_{i=1}^{216} 1-a_i=215$, we find $c=\frac{1}{215}$ and thus $x_2=c(1-a_2)=\frac{1}{215}\cdot \left(1-\frac{1}{4}\right)=\frac{3}{860}$. Hence the desired answer is $860+3=\boxed{863}$.

 





\textbf{Solution 2}

Squaring $\sum_{i=1}^{216}x_i$ yields\begin{gather*}\sum_{i=1}^{216}x_{i}^{2}+2\sum_{1\le i < j\le 216}x_{i}x_{j}=1\\\frac{214}{215}+\sum_{i=1}^{216}\frac{x_{i}^{2}}{1-a_{i}}=1\\\frac{4}{3}x_{2}^{2}=\frac{1}{215}-\sum_{i\neq 2}\frac{x_{i}^2}{1-a_{i}}\end{gather*}By the Cauchy-Schwarz Inequality,
 \[(\sum_{i\neq 2}\frac{x_{i}^2}{1-a_{i}})(\sum_{i\neq 2}1-a_{i})\ge (\sum_{i\neq 2}x_{i})^2\]The maximal of $x_{2}$ occurs as $\sum_{i\neq 2}\frac{x_{i}^2}{1-a_{i}}$, on the RHS, is minimized. We see that $\sum_{i\neq 2}1-a_{i}=\frac{857}{4}$. Thus,
 \[\frac{4}{3}x_{2}^{2}=\frac{1}{215}-\frac{4(1-x_{2})^2}{857}\]Solving yields $x_{2}=\frac{3}{860}$, and the requested sum is $3+860=\boxed{863}$.

 


 




\end{enumerate}